\documentclass[../RFC-HDFG-2026-001.tex]{subfiles}

\begin{document}

%% ======================================================================
%%  Section 10 – Compatibility
%% ======================================================================
\section{Compatibility}
\label{sec:compatibility}

\subsection{Semantic Change in \texttt{H5Z\_class3\_t}}
\label{sec:name-semantic-change}

Upgrading a filter plugin from \texttt{H5Z\_class2\_t} to \texttt{H5Z\_class3\_t}
introduces no source-level incompatibility at existing \emph{call sites}
(code that calls \texttt{H5Zregister} continues to compile and link without
modification). However, the plugin binary itself \emph{is} binary incompatible:
the plugin shared library must be recompiled against the new struct definition,
and any existing binary that exports an \texttt{H5Z\_class2\_t} cannot simply be
reused as an \texttt{H5Z\_class3\_t} plugin without recompilation. Plugin
distributors must ship updated binaries.

\texttt{H5Z\_class3\_t} keeps the \texttt{name} field at the same
offset and type as in \texttt{H5Z\_class2\_t}, but strengthens its
contract: in v3 \texttt{name} is required and must carry the filter's
canonical name (a stable, unique string identifier).  Filters are
still addressed by integer ID for programmatic dispatch (assigned by
the HDF Group filter registry~\cite{hdf5registeredfilters}), but
\texttt{name} now serves as both the canonical string identifier and
the human-readable label returned by \texttt{H5Pget\_filter2}.  The
library writes \texttt{name} into the filter-pipeline message's name
slot at filter-add time --- the same on-disk slot v2 plugins use ---
so the name persists in the file and is available to tools even when
the plugin is not installed.  No trailing \texttt{cd\_values} slots
are added.

\verysubsection{Before (v2): freeform debug comment}
In \texttt{H5Z\_class2\_t}, the \texttt{name} field is documented as a ``Comment for
debugging'' and is also user-visible through the \texttt{name[]} output parameter of
\texttt{H5Pget\_filter2}. Existing plugins commonly populate this field with
human-readable prose such as \texttt{"HDF5 ZFP filter, version 1.0.0 (rate,
precision, and accuracy modes)"} or \texttt{"My Compression Library v2.3"}.

\verysubsection{After (v3): canonical name in the pipeline-message name slot}
\texttt{H5Z\_class3\_t} keeps \texttt{name} but strengthens its
contract: it is now a required, stable, unique string identifier
(e.g., \texttt{"zfp"}, \texttt{"blosc2"}) that doubles as the
human-readable display label.  The library writes it into the
filter-pipeline message's name slot at filter-add time, which is the
same on-disk slot v2 plugins have always used.  \texttt{cd\_values}
is unchanged.

\verysubsection{Impact on \texttt{H5Pget\_filter2}}
For v3 plugins, \texttt{H5Pget\_filter2} returns the plugin's
\texttt{name} from the runtime registry when the plugin is loaded.
When the plugin is not loaded, it falls back to the name stored in
the filter-pipeline message's name slot at filter-add time.  If
neither is available, it falls back to the numeric filter ID
formatted as a decimal string.

\verysubsection{Migration guidance}

\begin{itemize}
\item \textbf{Pick a stable canonical name.}  When upgrading to v3,
  set \texttt{name} to a short, stable string identifier for the
  plugin class (e.g., \texttt{"zfp"}, \texttt{"blosc2"}).  Do not
  include version numbers, mode descriptors, or other variable
  content --- those belong in the configuration parameters surfaced by
  \texttt{H5Z\_get\_config\_func\_t}.  For example, a v2 plugin with
  \texttt{name = "HDF5 ZFP filter, version 1.0.0 (rate mode)"} should
  become a v3 plugin with \texttt{name = "zfp"}.
\item \textbf{No action required for consumers of \texttt{H5Pget\_filter2}
  that display the string verbatim:} These applications continue to
  function correctly; the displayed string is now the plugin's
  canonical name.
\end{itemize}

\subsection{On-Disk Format}

This RFC introduces no on-disk format changes.  The string configuration
is resolved at configuration time into \texttt{cd\_values} using the
existing \texttt{H5O\_PLINE\_VERSION\_2} format---output is byte-identical
to an equivalent \texttt{H5Pset\_filter} call.  Typed on-disk parameter
storage (per-slot type tags, which would enable exact round-trip of
floating-point parameters) is deferred to a follow-on RFC that would
extend \texttt{H5O\_PLINE\_VERSION}.

\subsection{Source Compatibility}

All existing public APIs are unchanged and not deprecated. Application code
using \texttt{H5Pset\_filter} continues to compile and function without modification.

\subsection{Binary / ABI Compatibility}

\texttt{H5Z\_class3\_t} is a new type, not a modification of \texttt{H5Z\_class2\_t}. New
functions are added to the public API. This is an API extension, not a
breaking change. A shared library minor version bump is appropriate.

\subsection{Plugin Compatibility}

Existing \texttt{H5Z\_class2\_t} plugins load and function without modification. They
can be used with the new API only via numeric ID with a raw \texttt{cd\_values}
parameter, or if they require no parameters. Plugin authors can upgrade to
\texttt{H5Z\_class3\_t} at their discretion.


%% ======================================================================
%%  Section 11 – Risks
%% ======================================================================
\section{Risks}
\label{sec:risks}

\verysubsection{Two-pass \texttt{set\_config} developer error}
The two-pass size-query protocol requires \texttt{set\_config} to report the
same \texttt{cd\_nelmts} on both invocations.  A developer who uses stateful
parsing (e.g., \texttt{strtok} on a modified copy of the string, or a
scan-forward pointer that is not reset) may produce a different count on the
second pass, triggering a library error or, in the worst case, a buffer
overrun before the library's second-pass capacity check catches it.
Mitigation: the library enforces the invariant by comparing first-pass and
second-pass \texttt{cd\_nelmts} and rejecting mismatches; the
\texttt{set\_config} template in \SectionRef{sec:set-config-template}
demonstrates the safe pattern; and the callback contract tests
(\texttt{cb-02}, \texttt{cb-03}) exercise mismatch detection.  Documentation
in \texttt{H5Zdevelop.h} explicitly warns against stateful parsing.

\verysubsection{Lossy round-trip user expectations}
Users and test suites often expect strict string equality after a
set-then-get round trip (\texttt{set\_config(s)} $\rightarrow$
\texttt{get\_config()} $\rightarrow$ $s' = s$).  Because
\texttt{cd\_values} is a lossy encoding for floating-point parameters
(e.g., a \texttt{double} rate quantised through \texttt{unsigned int}
slots), the reconstructed string $s'$ may differ from $s$ in decimal
representation even when the configuration is semantically identical.
The RFC addresses this in two places: the introspection note in
\SectionRef{sec:introspection} states that reconstructed strings are
informational only, and the round-trip tests (\SectionRef{sec:test-builtin-roundtrip})
verify \texttt{cd\_values} equality rather than string equality.  The
API documentation for \texttt{H5Pget\_filter\_params\_by\_idx} repeats
this caveat so that callers are not surprised.

\verysubsection{Grammar maintenance}
The TOML inline table subset defined in \SectionRef{sec:param-overview}
constitutes a specification the library must support indefinitely. Future
grammar extensions are backward-compatible: any string valid under the current
grammar remains valid with the same semantics under future grammars. The
semicolon is reserved in anticipation of a possible pipeline-string extension.

\verysubsection{Plugin adoption}
The value of the string API scales with the number of
filters that implement \texttt{set\_config}. If adoption is slow, v2 plugins can only
be used via the existing \texttt{H5Pset\_filter} API. Clear documentation and a low
barrier to v3 adoption mitigate this.

\verysubsection{Parsing overhead for small-data workloads}
The string API adds parsing cost at \texttt{H5Pappend\_filter} call time:
one linear pass over the parameter string (bounded by
\texttt{H5Z\_CONFIG\_STRING\_MAX} = 4096~bytes) per filter, followed by two
invocations of the \texttt{set\_config} callback (the two-pass size-query
protocol). For typical scientific workloads where dataset creation is
infrequent relative to I/O, this overhead is negligible.

In ``small data'' workloads that create thousands of lightweight datasets in a
tight loop, the per-dataset overhead of string parsing could become measurable,
though the dominant cost is already property list manipulation and object header
creation.

Applications that are sensitive to dataset-creation latency and configure the
same filter pipeline repeatedly should construct the DCPL once using
\texttt{H5Pappend\_filter} and reuse it across multiple \texttt{H5Dcreate}
calls, avoiding repeated parsing. This is already the recommended practice for
\texttt{H5Pset\_filter} and does not change with the new API.

\noindent\textbf{Decision:} The raw-string callback interface is retained.
\texttt{set\_config} receives the parameter string as \texttt{const char
*params} and uses \texttt{H5Zconfig\_get\_param} for key lookup.  Parsing
overhead is negligible: a typical filter string is under 100~bytes with
fewer than five keys, making two-pass parsing cost microseconds against a
millisecond dataset-creation operation.  Switching to a tokenized interface
would change the \texttt{H5Z\_set\_config\_func\_t} ABI, burden plugin
authors with a new token-array API, and foreclose future flexibility.
If parsing ever becomes measurable, token caching can be added inside the
library transparently without changing the callback interface.


%% ======================================================================
%%  Section 12 – Test Plan
%% ======================================================================
\section{Test Plan}
\label{sec:test-plan}

The following test categories are satisfied before the implementation is
merged.  Test names in \texttt{typewriter} font are illustrative identifiers
for the test matrix; the actual test suite may group them differently.

\subsection{Parser Unit Tests}
\label{sec:test-parser}

\noindent The parameter string parser (\SectionRef{sec:param-overview}) is the
lowest-level component and is exercised independently of any filter
registration.  Tests invoke the typed accessor functions
(\SectionRef{sec:config-get-param}) directly with hand-crafted strings.

\verysubsection{Valid inputs}

\begin{center}
\small
\renewcommand{\arraystretch}{1.25}
\begin{tabularx}{\textwidth}{l>{\raggedright\arraybackslash\hsize=0.7\hsize}X>{\raggedright\arraybackslash\hsize=1.3\hsize}X}
\toprule
\textbf{Test ID} & \textbf{Input} & \textbf{Expected result} \\
\midrule
\texttt{parse-01} & \texttt{level = 6}
  & \texttt{H5Zconfig\_get\_int} returns \texttt{6} \\
\texttt{parse-02} & \texttt{mode = "rate", rate = 3.5}
  & \texttt{H5Zconfig\_get\_str} returns \texttt{"rate"}; \texttt{H5Zconfig\_get\_double} returns \texttt{3.5} \\
\texttt{parse-03} & \texttt{fast\_mode = true, level = 6}
  & \texttt{H5Zconfig\_get\_bool} returns \texttt{TRUE}; \texttt{H5Zconfig\_get\_int} returns \texttt{6} \\
\texttt{parse-04} & \texttt{NULL}
  & valid; zero parameters; all accessors return \texttt{0} (absent) \\
\texttt{parse-05} & \texttt{""}
  & valid; zero parameters \\
\texttt{parse-06} & \texttt{path = "/data/run\_1,v2/dict.bin"}
  & comma inside quoted string is literal; \texttt{H5Zconfig\_get\_str} returns \texttt{/data/run\_1,v2/dict.bin} \\
\texttt{parse-07} & \texttt{note = "say \textbackslash"hi\textbackslash""}
  & TOML escape sequence; \texttt{H5Zconfig\_get\_str} returns \texttt{say "hi"} \\
\texttt{parse-08} & \texttt{\ \ level\ \ =\ \ 6\ \ }
  & whitespace stripped; \texttt{H5Zconfig\_get\_int} returns \texttt{6} \\
\texttt{parse-09} & \texttt{LEVEL = 6}
  & key normalized to lowercase \texttt{level}; \texttt{H5Zconfig\_get\_int} returns \texttt{6} \\
\texttt{parse-10} & \texttt{name = 'hello world'}
  & single-quoted literal string; \texttt{H5Zconfig\_get\_str} returns \texttt{hello world} \\
\texttt{parse-11} & \texttt{\{level = 6\}}
  & braced form accepted; result identical to bare form \\
\texttt{parse-12} & \texttt{stripe = \{size = 1048576, count = 4\}}
  & nested inline table; \texttt{H5Zconfig\_get\_int} with key \texttt{stripe.size} returns \texttt{1048576} \\
\texttt{parse-13} & \texttt{stripe.size = 1048576, stripe.count = 4}
  & dotted keys; equivalent to parse-12 \\
\texttt{parse-14} & \texttt{rate = -3.5}
  & negative float; \texttt{H5Zconfig\_get\_double} returns \texttt{-3.5} \\
\bottomrule
\end{tabularx}
\end{center}

\verysubsection{Malformed inputs (must return \texttt{H5E\_BADVALUE})}

\begin{center}
\small
\renewcommand{\arraystretch}{1.25}
\begin{tabularx}{\textwidth}{lX}
\toprule
\textbf{Test ID} & \textbf{Input and reason for rejection} \\
\midrule
\texttt{parse-E01} & \texttt{level = "rate"} passed to \texttt{H5Zconfig\_get\_int} --- type mismatch \\
\texttt{parse-E02} & \texttt{level = 6, level = 6} --- duplicate key \\
\texttt{parse-E03} & \texttt{path = "/no/closing/quote} --- unbalanced double-quote \\
\texttt{parse-E04} & String longer than \texttt{H5Z\_CONFIG\_STRING\_MAX} (4096) bytes --- exceeds maximum \\
\texttt{parse-E05} & String with more than \texttt{H5Z\_CONFIG\_MAX\_PARAMS} (64) top-level keys --- exceeds maximum \\
\texttt{parse-E06} & \texttt{rate = inf} --- special float value rejected by library \\
\texttt{parse-E07} & \texttt{rate = nan} --- special float value rejected by library \\
\texttt{parse-E08} & \texttt{data = [1, 2, 3]} --- arrays not in supported subset \\
\texttt{parse-E09} & bare semicolon outside quoted string --- reserved delimiter \\
\bottomrule
\end{tabularx}
\end{center}

\verysubsection{Boundary conditions}

\begin{itemize}
\item \texttt{parse-B01}: Single character key and integer value (\texttt{a = 1}).
\item \texttt{parse-B02}: Maximum allowed string length
  (\texttt{H5Z\_CONFIG\_STRING\_MAX} = 4096 bytes) --- accepted.
\item \texttt{parse-B03}: \texttt{H5Z\_CONFIG\_STRING\_MAX}+1 byte string --- rejected
  with \texttt{H5E\_BADVALUE}.
\item \texttt{parse-B04}: Exactly \texttt{H5Z\_CONFIG\_MAX\_PARAMS} (64) top-level keys ---
  accepted.
\item \texttt{parse-B05}: \texttt{H5Z\_CONFIG\_MAX\_PARAMS}+1 (65) top-level keys --- rejected
  with \texttt{H5E\_BADVALUE}.
\item \texttt{parse-B06}: Semicolon inside a double-quoted string value --- accepted.
\item \texttt{parse-B07}: Two-level nested inline table --- accepted; three-level --- accepted
  (depth is not bounded beyond the string length limit).
\end{itemize}


\subsection{Callback Contract Tests}
\label{sec:test-callbacks}

These tests verify the two-pass \texttt{set\_config} contract
(\SectionRef{sec:callback-specs}) and the \texttt{get\_config} round-trip guarantee.

\begin{itemize}
\item \texttt{cb-01}: First-pass call with \texttt{cd\_values=NULL} returns the correct
  \texttt{cd\_nelmts}; second-pass call populates \texttt{cd\_values} without exceeding
  the allocated buffer.
\item \texttt{cb-02}: A mock callback that returns a larger \texttt{cd\_nelmts} on the
  second pass triggers a library error (not a buffer overrun).
\item \texttt{cb-03}: A mock callback that returns a smaller \texttt{cd\_nelmts} on the
  second pass triggers a library error (not silent truncation).
\item \texttt{cb-04}: A callback that modifies \texttt{flags} does so silently;
  verify the updated flags value by comparing the input flags with those
  returned by \texttt{H5Pget\_filter2} after the call.
\item \texttt{cb-05}: A callback that does NOT modify \texttt{flags} leaves the
  error stack empty after the call.
\item \texttt{cb-06}: \texttt{get\_config} called with \texttt{buf=NULL} returns the required
  length without writing any bytes.
\item \texttt{cb-07}: \texttt{get\_config} output is a valid \texttt{set\_config} input (round-trip
  grammar check: parse the output and confirm it tokenizes without error).
\item \texttt{cb-08} (\textbf{modify-filter pattern}): Since no \texttt{H5Pmodify\_filter2}
  exists, users who need to update parameters of an existing filter in a copied
  DCPL must retrieve the current \texttt{cd\_values} via \texttt{H5Pget\_filter\_by\_id2}
  and call \texttt{H5Pmodify\_filter} directly.  Verify the pattern:
  \begin{enumerate}
    \item Create a DCPL and append deflate with \texttt{level=6} via
          \texttt{H5Pappend\_filter}.
    \item Copy the DCPL (\texttt{H5Pcopy}) to simulate a caller who received it.
    \item Call \texttt{H5Pget\_filter\_by\_id2} and confirm \texttt{cd\_values[0] == 6}.
    \item Set \texttt{cd\_values[0] = 9} and call \texttt{H5Pmodify\_filter}.
    \item Call \texttt{H5Pget\_filter\_by\_id2} again and confirm \texttt{cd\_values[0] == 9}.
  \end{enumerate}
\end{itemize}


\subsection{Built-in Filter Round-Trip Tests}
\label{sec:test-builtin-roundtrip}

For each built-in filter that implements \texttt{set\_config}/\texttt{get\_config}, a
round-trip test confirms that:
\begin{align*}
  &\text{set\_config}(s) \;\Rightarrow\; cd\_values_1
  \quad\text{and}\quad
  \text{get\_config}(cd\_values_1) \;\Rightarrow\; s' \\
  &\text{set\_config}(s') \;\Rightarrow\; cd\_values_2
  \quad\text{with}\quad cd\_values_1 = cd\_values_2
\end{align*}
The reconstructed string $s'$ need not equal $s$ literally; only the
resulting \texttt{cd\_values} must be identical.

\begin{center}
\small
\renewcommand{\arraystretch}{1.25}
\begin{tabularx}{\textwidth}{llX}
\toprule
\textbf{Test ID} & \textbf{Filter} & \textbf{Representative parameter string} \\
\midrule
\texttt{rt-01} & deflate    & \texttt{level = 0} (minimum) \\
\texttt{rt-02} & deflate    & \texttt{level = 6} (default) \\
\texttt{rt-03} & deflate    & \texttt{level = 9} (maximum) \\
\texttt{rt-04} & szip       & \texttt{coding = "entropy", pixels\_per\_block = 2} \\
\texttt{rt-05} & szip       & \texttt{coding = "nn", pixels\_per\_block = 32} \\
\texttt{rt-06} & scaleoffset & \texttt{scale\_type = "float\_dscale", scale\_factor = 4} \\
\texttt{rt-07} & scaleoffset & \texttt{scale\_type = "float\_escale", scale\_factor = 0} \\
\texttt{rt-08} & scaleoffset & \texttt{scale\_type = "int", scale\_factor = 8} \\
\texttt{rt-09} & shuffle    & \texttt{NULL} (parameterless) \\
\texttt{rt-10} & fletcher32 & \texttt{NULL} (parameterless) \\
\texttt{rt-11} & nbit       & \texttt{NULL} (parameterless) \\
\bottomrule
\end{tabularx}
\end{center}

\noindent Additional round-trip tests verify that:
\begin{itemize}
\item \texttt{rt-E01}: Passing a non-NULL, non-empty \texttt{params} to a parameterless
  built-in (\texttt{shuffle}, \texttt{fletcher32}, \texttt{nbit}) returns \texttt{H5E\_BADVALUE}.
\item \texttt{rt-E01b}: Passing \texttt{NULL} or an empty parameter string
  to a filter that implements \texttt{set\_config} and requires parameters
  (e.g., Scale-Offset, SZIP) invokes \texttt{set\_config} with
  \texttt{params~=~NULL} and returns \texttt{H5E\_BADVALUE} at
  \texttt{H5Pappend\_filter} time (fail-fast). The pipeline must not be
  appended.
\item \texttt{rt-E02}: Passing an unrecognized key (e.g., \texttt{clevl = 6}
  instead of \texttt{clevel = 6}) to any built-in returns
  \texttt{H5E\_BADVALUE} with an error message identifying the unrecognized key.
  v3 filters reject unknown keys (\SectionRef{sec:callback-specs}).
\item \texttt{rt-E03}: Passing an out-of-range value (\texttt{level = 10} for deflate,
  \texttt{pixels\_per\_block = 33} for szip) returns \texttt{H5E\_BADVALUE}.
\item \texttt{rt-E04}: Passing a type-mismatched value (\texttt{level = "six"} for deflate)
  returns \texttt{H5E\_BADVALUE} (TOML parser rejects the type before \texttt{set\_config}
  is invoked).
\end{itemize}


\subsection{Regression Tests: v1/v2 Plugin Behavior Unchanged}
\label{sec:test-regression}

These tests confirm that no existing behavior is broken.

\begin{itemize}
\item \texttt{reg-01}: A filter registered as \texttt{H5Z\_class2\_t} remains callable via
  \texttt{H5Pset\_filter} with the same integer ID and \texttt{cd\_values} as before.
  Data written and read through such a filter before and after the feature
  introduction is byte-identical.
\item \texttt{reg-02}: \texttt{H5Pappend\_filter} with \texttt{params=NULL} for a v2 filter
  succeeds (\texttt{NULL} is treated as ``no parameters'', equivalent to
  \texttt{H5Z\_PARAMS\_RAW(0, NULL)}).
\item \texttt{reg-03}: \texttt{H5Pappend\_filter} with \texttt{H5Z\_PARAMS\_STRING} for a
  v2 filter that does NOT implement \texttt{set\_config} returns \texttt{H5E\_UNSUPPORTED}.
\item \texttt{reg-04a}: The existing \texttt{H5Pget\_filter2} continues to return the
  \texttt{name} field from the class struct (the debug comment) for v2 plugins.
  Verify that a v2 plugin with \texttt{name = "HDF5 deflate filter v1.0"} returns
  that string.
\item \texttt{reg-04b}: For a v3 plugin, \texttt{H5Pget\_filter2}
  returns \texttt{name} from the runtime registry when the plugin is
  loaded.  For example, a plugin with \texttt{name = "zfp"} returns
  \texttt{"zfp"}.
\item \texttt{reg-04c}: For a v3 plugin that is not loaded but whose
  filter is present in a file, \texttt{H5Pget\_filter2} returns the
  canonical name stored in the filter-pipeline message's name slot at
  filter-add time.  If no stored name is available, it falls back to
  the numeric filter ID as a decimal string.
\item \texttt{reg-05}: All existing \texttt{h5repack} invocations with the legacy filter
  syntax (\texttt{GZIP}, \texttt{SHUF}, \texttt{FLET}, \texttt{UD=id,...}) produce
  output files byte-identical to those produced by the pre-feature version of
  \texttt{h5repack}.
\item \texttt{reg-06} (\textbf{cross-version read}): Write a dataset
  using a v3-registered plugin.  Read the file using a build of the
  library that presents only a v2 version of the same plugin.  Verify
  that the dataset is read correctly.  Because the v3 path writes
  the plugin's \texttt{name} into the pipeline-message name slot
  (already a v2 on-disk feature) and adds no trailing
  \texttt{cd\_values} slots, the v2 plugin sees exactly the
  \texttt{cd\_nelmts} it would have seen for an
  \texttt{H5Pset\_filter}-written dataset.
\item \texttt{reg-07} (\textbf{DCPL encode/decode round-trip}): Create a DCPL, append a
  v3 filter via \texttt{H5Pappend\_filter} with a string parameter, then
  \texttt{H5Pencode} the DCPL to a buffer.  \texttt{H5Pdecode} the buffer into a new
  DCPL.  Verify that the decoded DCPL contains identical filter ID, flags, and
  \texttt{cd\_values}---confirming that \texttt{set\_config} is not re-invoked on
  decode.
\end{itemize}


\subsection{Parallel HDF5 Tests}
\label{sec:test-parallel}

These tests require a Parallel HDF5 build and at least two MPI ranks.

\begin{itemize}
\item \texttt{par-01} (\textbf{DCPL consistency}): All ranks call
  \texttt{H5Pappend\_filter} with the same filter ID and parameter string; verify
  that the resulting \texttt{cd\_values} in the DCPL are byte-identical across ranks
  before \texttt{H5Dcreate}.
\item \texttt{par-02} (\textbf{rank failure propagation}): Simulate a missing plugin
  on one rank (e.g., by unsetting \texttt{HDF5\_PLUGIN\_PATH} on rank~1 only).
  Verify that the failing rank returns an error from \texttt{H5Pappend\_filter} (or from
  \texttt{H5Dcreate}, depending on when plugin loading occurs) and that the
  application can detect this and abort collectively without an MPI deadlock.
\item \texttt{par-03} (\textbf{debug-build consistency check}): In a debug build
  (\texttt{NDEBUG} not defined), induce a rank-inconsistent DCPL (by having one rank
  add an extra filter entry) and verify that \texttt{H5Dcreate} returns
  \texttt{H5E\_CANTCREATE} with an informative diagnostic rather than silently
  producing a corrupt file.
\item \texttt{par-04} (\textbf{built-in-only pipeline}): Create a chunked dataset using
  only built-in filters configured via the string API in a multi-rank parallel
  job; verify that the dataset can be written and read back correctly from all
  ranks.
\end{itemize}


\subsection{Thread-Safety Tests}
\label{sec:test-threadsafety}

These tests require a thread-safe HDF5 build (\texttt{--enable-threadsafe}).

\begin{itemize}
\item \texttt{ts-01} (\textbf{concurrent param lookups}): Spawn $N$ threads
  ($N \geq 8$) that each call \texttt{H5Pget\_filter\_params\_by\_idx} for built-in filter
  pipeline entries in a tight loop.  No data race or crash occurs.
\item \texttt{ts-02} (\textbf{concurrent registration and param lookup}): One thread
  calls \texttt{H5Zregister} with a new v3 plugin while other threads concurrently
  call \texttt{H5Pget\_filter\_params\_by\_idx} on pipeline entries that may include that
  filter.  Lookups see a consistent state (plugin absent or present);
  no partial or corrupt state.
\item \texttt{ts-03} (\textbf{concurrent \texttt{H5Pappend\_filter}}): Threads each create
  their own DCPL and call \texttt{H5Pappend\_filter} concurrently for the same built-in
  filter ID.  All calls succeed without error or data corruption (property
  lists are per-object and not shared).
\item \texttt{ts-04} (\textbf{on-demand plugin scan under contention}): Arrange for a
  plugin to be loaded on-demand (not yet registered).  Then have $N$ threads
  simultaneously call \texttt{H5Zfilter\_avail} for that filter's ID.  Verify the
  plugin is registered exactly once and all threads receive a positive result.
\end{itemize}

\end{document}
